\section{Observation noise and its removal before the pipeline}
\label{app:noise}

\Cref{def:m1} and \cref{thm:ng} describe a noiseless image $X^u=g(Z^u)$. An image is
\emph{environment-specific} even without noise: its latents $Z^u$ have a law that depends on the
environment $u=(y,t,e)$ (content through $c(u)=(y,t)$, style through $r(u)=(t,e)$), so the clean image
$g(Z^u)$ is already location-specific. The only objects assumed to be shared across environments are the
mixing $g$ and the noiseless mechanism it defines (\labelcref{item:P1,item:P2}). This appendix extends
the model by an additive observation noise whose level also depends on the environment (the camera and
its location), and asks what happens downstream and what to do about it. The extension is
\begin{equation}
X^u \;=\; g(Z^u)+\eta^u,\qquad \eta^u\sim\Nn\bigl(0,\varsigma_u^2 I_d\bigr),
\label{eq:noise-obs}
\end{equation}
with the family (Gaussian) assumed, the level $\varsigma_u$ unknown, and $\eta^u$ independent of $Z^u$
given $u$. In M1 the level is allowed to depend on $u$ through $r(u)=(t,e)$ (camera, and day or night
mode), not through the species $y$; this is an assumption of this appendix, not of \cite{ng2025}.
The proposal is to \emph{clean} the image with a fixed, deterministic operation before it enters the
frozen backbone of \cref{sec:pseudoinj:invae}. Nothing here has been run. The arguments about
propagation are first-order arguments of mine, not results from \cite{ng2025} or any other paper, and
are marked as such.

\subsection{What is shared and what is environment-specific}
\label{app:noise:inv}
\paragraph{Convention.} A symbol carries the superscript $u$ (or $(u)$ for a density) \ifft\ 
its value or its law depends on the environment. Maps that are the same for every environment carry no
index. Realizations drawn from environment $u$ keep the index: $z^u$, $s^u$, $x^u$.
In particular the clean image is never written as a location-free $g(z)$: it is $S^u=g(Z^u)$.

\begin{center}
\small
\begin{tabular}{@{}L{5.2cm} L{9.8cm}@{}}
\toprule
\multicolumn{2}{@{}l}{\textbf{shared across environments}}\\
\midrule
mixing $g$ & \labelcref{item:P1}; also the manifold $\mathcal M=g(\R^n)$ it spans (the manifold is
shared, the law on it is not)\\
noiseless mechanism: \newline
$p(s\mid z)=\delta_{g(z)}$ & the clean image given the latents; the same Dirac at
every location. This is the only mechanism assumed invariant.\\
DAG $\Gc$, content laws $p(z_c\mid c)$ & \labelcref{item:P2}; depend on $c=(y,t)$, not on $e$\\
backbone $h$, a fixed cleaning $\mathrm{Cl}$ & design choices applied identically at every
location and at inference\\
\midrule
\multicolumn{2}{@{}l}{\textbf{specific to an environment}}\\
\midrule
latent law $p^{(u)}_Z$ & content through $c(u)$, style through $r(u)$\\
clean image $S^u=g(Z^u)$ and its law $p^{(u)}_S$ & the map is shared, the image is not: camera,
geography and orientation enter through $Z_s^u$\\
observation noise $\eta^u$, level $\varsigma_u$ & new in this appendix\\
observed image $X^u$, law $p^{(u)}_X$ & $p^{(u)}_X=p^{(u)}_S*\Nn(0,\varsigma_u^2I)$\\
observed conditional: \newline
$p^{(u)}(x\mid z)=\Nn(g(z),\varsigma_u^2I)$ & \emph{not} invariant once the noise is
location-dependent; only its mean map $g$ is\\
cleaned image, bias, residual, cached $\tilde x^u$, estimates $\hat z^u$ & all derived from the above\\
\bottomrule
\end{tabular}
\end{center}

\paragraph{Consequence for what cleaning must do.} The task is to remove $\eta^u$ and to keep the
location-specificity of $S^u$. Cleaning is not meant to make images location-free: the style latents
$Z_s^u$ live in $S^u$ and the style block needs them. The line between noise and style is therefore
this: whatever is generated through $g$ from $Z^u$ is signal (style included); whatever is added after
$g$, independently of $Z^u$, is noise. Camera grain sits on this line. Grain that is a function of the
style latents is style; grain that is fresh independent randomness per pixel is noise. Real grain is
probably a mixture, and no cleaning step can separate the two exactly.

\subsection{Terminology}
\label{app:noise:terms}
The words below are used in one sense only in this appendix. Earlier informal use of ``deconvolving an
image'' is replaced by ``denoising'' (image level) or ``deconvolution'' (law level), see the table.

\begin{center}
\small
\begin{tabular}{@{}l L{3.4cm} L{1.2cm} L{6.6cm}@{}}
\toprule
term & symbol & depends on $u$? & meaning\\
\midrule
signal (clean image) & $S^u=g(Z^u)\in\R^d$ & yes & the clean image of environment $u$, pixel space; the
map $g$ is shared, the image is not\\
structural noise & $\varepsilon$ & yes & the noise inside the latent equations
\cref{eq:content,eq:style}; \emph{not} observation noise\\
observation noise & $\eta^u$, level $\varsigma_u$ & yes & additive noise of \cref{eq:noise-obs}, added to
the image after $g$\\
observed image & $X^u=S^u+\eta^u$ & yes & what the camera delivers\\
noise-level estimation & $\hat\varsigma$ & -- & estimating $\varsigma_u$ from images (per image, per cell,
or per location)\\
denoising / cleaning & $\mathrm{Cl}:\R^d\to\R^d$ & no, if fixed & a map acting on \emph{one image} that
returns an estimate of $S^u$ (not of a location-free image); \emph{image level}\\
deconvolution & -- & -- & undoing the convolution of laws,\newline
$p^{(u)}_X=p^{(u)}_S*\Nn(0,\varsigma_u^2I)$; \emph{law level}. It returns a distribution, never a clean
image\\
filtering & $F$ & no & denoising by a fixed operator (blur, block averaging, median, fixed wavelet
threshold)\\
fixed / adaptive & $\mathrm{Cl}$ / $\mathrm{Cl}_{\hat\varsigma}$ & no / yes & fixed: the same map for
every image and location; adaptive: its strength depends on $\hat\varsigma$, hence on the environment\\
equalization & -- & -- & adding noise so that every environment ends at a common level\\
projection & -- & -- & mapping $x$ to a nearest point of the shared manifold $\mathcal M=\{g(z)\}$
(ideal denoising)\\
bias of cleaning & $b^u(z)$ & yes, if $\mathrm{Cl}$ adapts & systematic distortion of the signal by
$\mathrm{Cl}$, at fixed $z$\\
residual noise & $r^u$ & yes & zero-mean fluctuation left after $\mathrm{Cl}$\\
\bottomrule
\end{tabular}
\end{center}
The bias and the residual are defined for the cleaned image $X'^u=\mathrm{Cl}(X^u)$ by
\begin{equation}
X'^u = g(Z^u)+b^u(Z^u)+r^u,\quad b^u(z)=\mathbb E[X'^u\mid Z^u=z]-g(z),\quad
r^u=X'^u-\mathbb E[X'^u\mid Z^u].
\label{eq:cleaned-model}
\end{equation}
Both are defined at fixed $z$, so any dependence of $b^u$ on $u$ comes from the cleaning (through the
noise level), not from the law of $Z^u$. Ideal cleaning means $b^u\equiv0$ and $r^u\equiv0$.
\newpage
\subsection{Where the noise enters, and how it propagates}
\label{app:noise:prop}
The stages, in order (shared maps are marked, everything else is environment-specific):
\begin{enumerate}[label=(\arabic*),leftmargin=2.2em]
\item\label{item:stage1} latent $Z^u$ (content and style), with the ANM prior of \cref{sec:use}, its law
      depending on $u$ through $c(u)$ and $r(u)$;
\item\label{item:stage2} clean image $S^u=g(Z^u)$, by the shared map $g$ (pixel space,
      $d\approx2\times10^5$ at $256\times256\times3$);
\item\label{item:stage3} \textbf{noise enters here:} $X^u=S^u+\eta^u$;
\item\label{item:stage4} optional cleaning $X'^u=\mathrm{Cl}(X^u)$ (this appendix), shared if fixed;
\item\label{item:stage5} frozen, shared IN-VAE encoder $h$ (the backbone), giving the cached
      $\tilde x^u=h(x'^u)$ of $32\times16\times16$ values;
\item\label{item:stage6} trainable head: $\tilde x^u\to\hat z^u\to\hat{\tilde x}^u$, with reconstruction
      and KL computed on $\tilde x^u$ only (\cref{sec:pseudoinj:arch});
\item\label{item:stage7} M1 prior and Ng machinery on $\hat z^u$.
\end{enumerate}
The head never sees pixels and has no pixel-space decoder.
Whatever is in $X'^u$ at stage \labelcref{item:stage4} is all
that stages \labelcref{item:stage5,item:stage6,item:stage7} can use.

\paragraph{Physical order versus the model.} A camera adds noise at native resolution, the file is then
JPEG-compressed, and the dataset images are resized to $256\times256$ in preprocessing
(\texttt{ti\_dataset.py} reads already resized images). So \cref{eq:noise-obs} at $256$ pixels is an
idealization: the effective noise is reduced by the resize, spatially correlated, and not exactly
independent of the signal. Only the level \emph{after} these steps matters, and only its dependence on
the environment.

\paragraph{No cleaning ($X'^u=X^u$).} First order (my argument): with $\Phi=h\circ g$ (both shared),
\[
\tilde x^u=h(g(z^u)+\eta^u)\approx\Phi(z^u)+J_h(g(z^u))\,\eta^u .
\]
The perturbation $J_h\eta^u$ splits into a part tangent to the manifold $\Phi(\R^n)$ and a normal part.
\begin{itemize}[leftmargin=1.5em]
\item \emph{Tangent part.} It is indistinguishable from a shift of $z^u$, so what is estimated is
      $z^u+\delta^u$, not $z^u$. The shift is correlated across coordinates and depends on $z^u$ through
      $J_h$. Consequences for the model: spurious edges in the Markov network (including between
      style and content, against \labelcref{item:P3}), parent-dependent residuals in the content
      equations (the ANM is violated), and, because $\varsigma_u$ depends on the environment, an
      estimated content law that depends on the location although $p(z_c\mid c)$ does not
      (against \labelcref{item:P2}).
\item \emph{Normal part.} It moves $\tilde x^u$ off the manifold. The head's reconstruction term treats
      it as unexplained residual. It does not enter $\hat z^u$ directly, but it competes with the
      latents for capacity.
\end{itemize}
Because $h$ is nonlinear and compresses $196{,}608$ values to $8{,}192$, the effective noise at the head
is neither additive, nor Gaussian, nor isotropic. A Gaussian likelihood with a per-environment
variance cannot be placed in the head. This is why the noise has to be handled before $h$.

\paragraph{With cleaning.} By \cref{eq:cleaned-model} the head sees $g(z^u)+b^u(z^u)+r^u$, and the two
errors act differently.
\begin{itemize}[leftmargin=1.5em]
\item \emph{Residual $r^u$.} It is additive and roughly independent of $z^u$. In pixel space the part of
      isotropic noise tangent to $\mathcal M$ has total variance $n\varsigma^2$, a fraction $n/d$ of the
      total, so it is the small part; the rest is off-manifold. What remains is the same kind of
      perturbation as without cleaning but much smaller. It is environment-specific if $\varsigma_u$ is
      and the cleaning does not equalize it.
\item \emph{Bias $b^u(z)$.} It changes the mixing: the effective map is $g_u=g+b^u$. If $b^u$ does not
      depend on $u$, then $g+b$ is again one shared map, and the only cost is some loss of
      injectivity. If $b^u$ depends on $u$, the mixing depends on the location, which is the
      camera-specific transformation listed as outside M1 and outside \cite{ng2025}
      (\cref{sec:limits}); content latents then move with the location. This is the one place where
      cleaning can break the shared mixing, which is the only thing assumed shared.
\end{itemize}

\paragraph{What is exact and what is not.}
Under \cref{eq:noise-obs} with known $\varsigma_u$, deconvolution at the law level is exact: the
characteristic function of a Gaussian never vanishes, so $p^{(u)}_X$ determines $p^{(u)}_S$ uniquely, and
\cref{thm:ng} then applies to the noiseless law as before. This says nothing about images: from a
single noisy image $x^u$ the clean image $s^u$ is not recoverable, only estimable through a prior, with
error that depends on that prior. So exact identifiability is a statement about the ideal case
$b^u\equiv r^u\equiv0$.
Identifiability does not hold exactly in practice anyway (pseudo-injectivity of $h$ and of the head,
finite samples, estimation error), so the useful question is when it holds \emph{nearly}. Working
criterion (not a theorem):

\noindent\hspace*{1cm}\fbox{\parbox{0.8\textwidth}{
cleaning is good enough when
the environment-dependent part of $b^u$ and the tangent part of $r^u$,
measured as a shift of $\hat\mu_c$, is small against the error floor of the noiseless pipeline,
where the floor is the variation of $\hat\mu$ between repeated runs on clean data (seeds, resamples).
}}

It must also stay below the margins of the method: the gaps between environments needed for sufficient
change, the margin used to call a node a sink, and the separation of species within a cell.
\vspace{\parskip}
\begin{remark}[Identifying $\varsigma_u$, a sketch of mine]
If the noise is isotropic Gaussian and $d\gg n$, the level should be identifiable at the law level: the
law of $X^u$ deconvolved by a trial level $\varsigma'\le\varsigma_u$ stays a valid law, and it is
singular (concentrated near an $n$-dimensional set) only at $\varsigma'=\varsigma_u$. This needs
isotropy and $d>n$; a general unknown covariance is confounded with the tangent directions. It is not
verified, and real noise here is not isotropic (see above).
\end{remark}

\subsection{Cleaning approaches}
\label{app:noise:methods}

\paragraph{0. No cleaning.} The baseline. All effects of \cref{app:noise:prop} are present.

\paragraph{1. Fixed filter.} One operator $F$ (Gaussian or box blur, block averaging, median, fixed
wavelet threshold), applied identically at every location and at inference. For a linear $F$,
\[
X'^u=Fg(Z^u)+F\eta^u,\qquad b=(F-I)g(z)\ \text{(independent of }u),\qquad r^u=F\eta^u .
\]
The map $F\circ g$ is shared across locations; the cleaned images are still location-specific through
$Z^u$, and the environment-specific part of the error sits only in the residual. Averaging $k$ pixels
reduces white noise variance by about $1/k$. Costs: it discards the
frequency band where the noise lives, and with it fur texture, edges and camera grain, which is part of
the style signal; it removes nothing inside the signal band; and the $1/k$ rule is optimistic for
resized, JPEG-compressed, signal-dependent noise. Choice of strength: too weak leaves
environment-specific residual noise, too strong loses information uniformly (less injective, but no
location leak). Bias toward over-filtering, with the noisiest training location as the reference.

\paragraph{2. Adaptive filter.} The same family with strength set by $\hat\varsigma$:
$X'^u=F_{\hat\varsigma(u)}(g(Z^u)+\eta^u)$. In the ideal case (projection onto $\mathcal M$ at every
level) it is exact. A real adaptive filter removes noise by distorting the signal, and distorts it
differently per location: the bias $b^u$ depends on $u$ (the same animal, at the same $z$, is crisp at a
clean site and blurred at a noisy one), so the shared mixing is lost. Whether this is worse than the
over- or under-filtering of a fixed filter is a practical question, not a theoretical one. It can win
when noise levels differ strongly between sites, so that a filter tuned for the noisiest site
over-smooths the clean ones.

Estimating $\hat\varsigma$: from a robust high-frequency statistic (wavelet median absolute deviation,
flat-patch variance), computed on the background only, with the animal bounding boxes masked (they are
available in \texttt{ti\_dataset.py}), so that fur texture does not inflate it. The estimate is biased
for correlated, signal-dependent noise; only its consistency across locations matters. Use pooled
estimates per cell in training (many images) and the same estimator per image at inference.

\paragraph{3. Equalization.} Estimate $\hat\varsigma_u$ and add $\Nn\bigl(0,\varsigma_{\max}^2-\hat\varsigma_u^2\bigr)$
to every image, so that all environments end at the level $\varsigma_{\max}$. The noise stops depending on $u$, nothing is
removed from the signal, and there is no bias, but information is degraded to the noisiest site and the
common noise still passes through $h$. It is not deterministic: draw the added noise once per image with
a fixed seed, so that the cached features are reproducible. It equalizes exactly only for
signal-independent noise, and a mis-estimated $\hat\varsigma_u$ leaves an environment-specific residual of
the same size as the error. Note that \emph{subtracting} a sampled noise draw does not denoise: the
real $\eta^u$ is a specific realization, an independent draw $\eta'$ leaves $\eta^u-\eta'$ with variance
$2\varsigma_u^2$.

\paragraph{4. Learned denoiser.} A denoising autoencoder (or a Noise2Noise or Noise2Void style network)
trained on the noisy images. At its optimum it returns the posterior mean
$x+\varsigma_u^2\nabla_x\log p^{(u)}_X(x)$ for Gaussian noise (Tweedie), which for small $\varsigma_u$ is close
to a projection onto the learned manifold, and that is the principled version of ideal denoising. But the manifold has to be
learned from the same data, a network that is trained jointly with $\hat\varsigma$ is circular, its
shrinkage again depends on $u$, and it may remove style. A generic pretrained denoiser projects onto the
natural-image manifold, which has a much higher dimension than $\mathcal M$, so more noise stays and some signal
goes. Not recommended as a first step.

\begin{center}
\small
\begin{tabular}{@{}l L{4.2cm} L{3.8cm} L{4.3cm}@{}}
\toprule
approach & where environment dependence lands & removes & main risk\\
\midrule
none & noise, through $h$ & nothing & all effects of \cref{app:noise:prop}\\
fixed filter & residual $r^u$ only & noise in one band & loses signal in the band, including grain\\
adaptive filter & bias $b^u$ and residual & noise in a band, per level & location-dependent mixing\\
equalization & residual, at a common level & nothing (adds) & information loss, reliance on $\hat\varsigma_u$\\
learned denoiser & bias $b^u$ and residual & noise off the learned manifold & circularity, style removal\\
\bottomrule
\end{tabular}
\end{center}

\subsection{A single test image}
\label{app:noise:single}
A test image $x^{u^*}$ comes from an unseen environment $u^*$, and nothing about $u^*$ is available
except the image itself. One image is enough to estimate the noise level: it has $10^5$ to $10^6$
pixels, so for roughly white
noise $\hat\varsigma$ is accurate, and per-image blind denoising is the standard setting of classical
denoisers. The practical limit is the bias of the estimator for correlated, signal-dependent noise, which
matters only if it differs by location. The rule is to use the \emph{same} per-image procedure in
training and at inference. Estimating per location in training and per image at test would give two
different pipelines. A fixed filter needs no estimate at all. Self-supervised single-image denoisers
(Noise2Self, Deep Image Prior) need no noise model but are slow, since they run on every image.

\subsection{Recommendation and test protocol}
\label{app:noise:rec}
\begin{enumerate}[label=(\arabic*),leftmargin=2.2em]
\item \textbf{Measure first.} Estimate $\hat\varsigma$ on the background of the images of each training
      location. If the levels are close, there is nothing to adapt, and a fixed filter (or none) is
      enough.
\item \textbf{Start with a fixed filter,} one strength for all locations, chosen for the noisiest
      training location, erring toward over-filtering. Reasons: it keeps the shared map $F\circ g$ shared, has one
      hyperparameter, and needs no extrapolation to a new location.
\item \textbf{Adaptive only if needed:} if the levels differ strongly and the fixed filter visibly
      over-smooths the clean sites. Then use the per-image estimator of \cref{app:noise:single}.
\item \textbf{Equalization} is the alternative to adaptivity when the signal-independent
      approximation is acceptable. A learned denoiser is the last resort.
\end{enumerate}
Decide between the arms (none, fixed, adaptive, optionally equalization) empirically:
\begin{itemize}[leftmargin=1.5em]
\item leave-one-location-out macro accuracy over L38/L43/L46 (L100 stays untouched until the end);
\item the noiseless error floor of $\hat\mu$ (repeated runs on clean data) and the shift of $\hat\mu_c$
      when known, location-dependent noise is added to the least noisy images before the arm is applied;
      the threshold of \cref{app:noise:prop} is where that shift crosses the floor;
\item the linear probe of \cref{app:probe} after cleaning (species accuracy, and within-species location
      accuracy, which must stay above chance), since cleaning must keep the location-specific style
      signal of $S^u$ while removing $\eta^u$.
\end{itemize}

\paragraph{Implementation notes.}
\begin{itemize}[leftmargin=1.5em]
\item Apply the same operator in training and at inference, before \texttt{extract\_features.py}, and
      keep the values in the range the extractor checks (clip after filtering).
\item \texttt{ti\_dataset.py} reads images that are already resized. A filter inserted there acts after
      the resize, on noise that is already averaged and correlated. Filtering before the resize means
      changing the preprocessing; the bounding-box transform is not affected, since only intensities change.
\item Brightness (the time-of-day proxy of \cref{sec:use}) is a mean and is nearly unaffected by
      zero-mean noise or low-pass filtering. Decide once whether it is computed before or after
      cleaning, and use that choice at every location.
\end{itemize}

\subsection{Open items}
\label{app:noise:open}
\begin{itemize}[leftmargin=1.5em]
\item The assumption that the noise level depends on $u$ only through $r(u)=(t,e)$, and that $\eta^u$ is
      independent of $Z^u$ given $u$, is untested. Grain that depends on the style latents violates the
      second part.
\item The tangent shift of $\hat z^u$ under residual noise is a first-order argument only; its size for the
      actual $h$ and head is not measured.
\item The sketch of identifying $\varsigma_u$ needs isotropy and $d>n$, neither of which holds for real
      feature residuals or for camera noise.
\item Which fixed filter family preserves style best, and whether $F\circ g$ stays close enough to
      injective, are untested; the linear probe is a necessary check, not a proof.
\item Whether the noise level differs enough between L38, L43 and L46 for adaptivity to matter has not
      been measured.
\end{itemize}